\hypertarget{introduction}{\section{Introduction}}
\label{sec:introduction}

% \tldr{Benchmarks are important, LMs are a thing)}
Benchmarks orient AI.
They encode values and priorities \citep{ethayarajh2020utility, birhane2022values} that specify directions for the 
AI community to improve upon \citep{jones1995evaluating, jones2005acl, kiela2021dynabench, bowman2021fix,  raji2021benchmark}.
When implemented and interpreted appropriately, they enable the broader community to better understand AI technology and influence its trajectory.

In recent years, the AI technology that has arguably advanced the most is foundation models \citep{bommasani2021opportunities}, headlined by the rise of language models \citep[LMs;][]{peters2018elmo, devlin2019bert,brown2020gpt3,rae2021gopher,chowdhery2022palm}.
At its core, a language model is a box that takes in text and generates text (\reffig{languageModel}).
Despite their simplicity,
when these models are trained on broad data
at immense scale, they can be adapted (\eg prompted or fine-tuned) to myriad downstream scenarios.
Yet the immense surface of model capabilities, limitations, and risks remains poorly understood.
The rapid development, rising impact, and inadequate understanding demand that we benchmark language models holistically.

\FigTop{figures/language_model_helm.png}{0.2}{languageModel}{
\textbf{Language model.}
A language model
takes text (a prompt) and generates text (a completion) probabilistically.
Despite their simple interface, language models can be adapted to a wide range of language tasks from
question answering to summarization.
}

% \tldr{What does it mean to holistically evaluate a language model}
But what does it mean to benchmark language models holistically?
Language models are general-purpose text interfaces that could be applied across a vast expanse of scenarios.
And for each scenario, we may have a broad set of desiderata: models should be accurate, robust, fair, efficient, and so on.
In fact, the relative importance of these desiderata often will depend not only on the perspective and values one has, but the scenario itself (\eg inference efficiency might be of greater importance in mobile applications).

\noindent We believe holistic evaluation involves three elements:
\begin{enumerate}
\item \textbf{Broad coverage and recognition of incompleteness.}
Given language models' vast surface of capabilities and risks,
we need to evaluate language models over a broad range of scenarios.
Broadening the evaluation has been a continuing trend in the NLP community, 
going from individual datasets such as SQuAD \citep{rajpurkar2016squad}
to small collections of datasets such as SuperGLUE \citep{wang2019glue}
to large collections of datasets such as the GPT-3 evaluation suite \citep{brown2020gpt3}, Eleuther AI LM Harness \citep{gao2021harness}, and BIG-Bench \citep{srivastava2022bigbench}.
However, it is neither possible to consider all the scenarios nor all the desiderata that (could) pertain to LMs.
Therefore, holistic evaluation should provide a top-down taxonomy and make explicit all the major scenarios and metrics that are missing.
\item \textbf{Multi-metric measurement.}
Societally beneficial systems reflect many values, not just accuracy.
Holistic evaluation should represent these plural desiderata, evaluating every desideratum for each scenario considered.
\item \textbf{Standardization.}
Our \emph{object} of evaluation is the language model, not a scenario-specific system.  Therefore, in order to meaningfully compare different LMs, the strategy for adapting an LM to a scenario should be controlled for.  Furthermore, each LM should be evaluated on the same scenarios to the extent possible.
\end{enumerate}
Overall, holistic evaluation builds transparency by assessing language models in their totality.
Rather than honing in on a specific aspect, we strive for a fuller characterization of language models to improve scientific understanding and orient societal impact.

% 1. Taxonomy of scenarios x metrics (aspirational)
% 2. Deliberate selection (given resources, access, best we can do)
\hypertarget{helm}{\subsection{\benchmarkname}}
\label{sec:helm}
Holistic Evaluation of Language Models (\benchmarkname) has two levels:
(i) an abstract taxonomy of \textit{scenarios} and \textit{metrics} to define the design space for language model evaluation 
and (ii) a concrete set of implemented scenarios and metrics that were selected to prioritize coverage (\eg different English varieties), value (\eg user-facing applications), and feasibility (\eg limited engineering resources).

\FigTop{figures/fig_taxonomy_v2}{0.35}{taxonomy}
{\textbf{The importance of the taxonomy to \benchmarkname.}
Previous language model benchmarks (\eg SuperGLUE, EleutherAI LM Evaluation Harness, BIG-Bench) are collections of datasets, each with a standard task framing and canonical metric, usually accuracy (\textit{left}).
In comparison, in \benchmarkname we take a top-down approach of first explicitly stating what we want to evaluate (\ie scenarios and metrics) by working through their underlying structure.
Given this stated taxonomy, we make deliberate decisions on what subset we implement and evaluate, which makes explicit what we miss (\eg coverage of languages beyond English).}

% Importance of the taxonomy
\paragraph{Recognition of incompleteness.}
Benchmarks across AI, including those for language models like SuperGLUE \citep{wang2019superglue}, the EleutherAI LM Harness \citep{gao2021harness}, and BIG-bench \citep{srivastava2022bigbench}, are defined by specific choices of scenarios and metrics.
Different benchmarks make different decisions on what to prioritize, how to make these decisions, and to what extent these processes are made clear in presenting the benchmark.
Since our aim is holistic evaluation, we believe it is necessary to be explicit on the relationship between what we aspire to evaluate and what we actually evaluate.
The construction of \benchmarkname starts top-down with a taxonomy over scenarios and metrics (see \reffig{taxonomy}).
The taxonomy not only facilitates the systematic selection of scenarios and metrics, but it also make explicit what is missing.
We view \benchmarkname as a living benchmark,
and we hope that both the abstract taxonomy and the concrete selection of scenarios and metrics will evolve according to the technology, applications, and social concerns.
In \refsec{missing}, we explicitly highlight evaluations \benchmarkname lacks that should be prioritized.  Often these are ones the entire AI field has historically neglected.

% What is the concrete benchmark
\FigTop{figures/fig_multi_metric_v2}{0.42}{multi-metric}
{\textbf{Many metrics for each use case.}
In comparison to most prior benchmarks of language technologies, which primarily center accuracy and often relegate other desiderata to their own bespoke datasets (if at all), in \benchmarkname we take a multi-metric approach.
This foregrounds metrics beyond accuracy and allows one to study the tradeoffs between the metrics.
}

\paragraph{Multi-metric measurement.}
\benchmarkname currently implements a \core\footnote{We use the term \textit{\core} to indicate that for this set of scenarios, we measure a range of metrics/desiderata. The term \textit{\core} is not meant to suggest that any specific scenario in this set is more fundamental than scenarios outside the set.} set of \numcorescenarios scenarios and \nummetriccategories (categories of) metrics.
Our scenarios, which are triples of (task, domain, language), span \numcoretasks user-facing tasks (\eg question answering, information retrieval, summarization, toxicity detection), several domains (\eg news, books), and currently only English (though we cover several English varieties such as African-American English and the English varieties spoken in different English-speaking countries). 
And our \nummetriccategories categories of metrics reflect a range of societal considerations (\ie accuracy, calibration, robustness, fairness, bias, toxicity, efficiency).
We emphasize that while we have specific quantitative metrics for all of these considerations, they (\eg fairness) are complex and contested social constructs that can be operationalized in many different ways. 
Consistent with our second element of holistic evaluation, 
we ensure our benchmark attains \textbf{dense} multi-metric measurement: of the 112 possible (\core scenario, metric) pairs, we measure 98 (87.5\%) as shown in \reftab{scenarios-metrics-matrix}.

This multi-metric perspective conveys a position we take on evaluation practices in AI. 
While most benchmarks primarily foreground accuracy, perhaps deferring the evaluation of other metrics (\eg the extent to which models generate toxic content) to separate scenarios (\eg \realtoxicityprompts), we believe it is integral that all of these metrics be evaluated in the same contexts where we expect to deploy models (see \reffig{multi-metric}).
In particular, measuring these \nummetriccategories desiderata for the same scenarios makes explicit potential trade-offs and helps to ensure these desiderata are not treated as second-class citizens to accuracy \citep[see][]{friedman1996bias}.

\paragraph{Targeted evaluations.}
In addition to our \core set of \numcorescenarios scenarios, where for each scenario we measure all \nummetriccategories categories of metrics, \benchmarkname has \numcomponents targeted evaluations through \numcomponentscenarios additional scenarios and accompanying metrics.
These evaluations target linguistic understanding, world and commonsense knowledge, reasoning capabilities, memorization and copyright, disinformation generation, biases, and toxicity generation, providing a deeper dive beyond the core scenarios.
This includes \numnewscenarios scenarios that are either entirely new (\eg \wikifact) or that have not been used in mainstream language model evaluation (\eg \ice).
%\pl{it might be nice to cite the various benchmarks when they are mentioned at least at first}
While \benchmarkname is oriented by a holistic approach that foregrounds societal impact and is reflected in our multi-metric perspective, evaluation can also pinpoint specific phenomena to advance scientific understanding \citep[\eg a model's ability to perform analogical reasoning; see][\S4.4]{bommasani2021opportunities}. 
For this reason, to make our evaluation results more intelligible, we separate the \core scenarios from the targeted evaluations: the \core scenarios and multi-metric measurement provide an integrated lens on models, whereas the targeted evaluations isolate specific skills and risks.

\FigTop{figures/fig_model_comparison_v4}{0.53}{evaluation-status-quo}
{\textbf{Standardizing language model evaluation.} Prior to our effort (\textit{top}), the evaluation of language models was uneven.
Several of our \numcorescenarios \core scenarios had no models evaluated on them, and only a few scenarios (\eg \boolq, \hellaswag) had a considerable number of models evaluated on them. 
Note that this is \textit{cumulative}: in the \textit{top} plot, we not only document instances where the work introducing the model evaluated on a given scenario, but \textbf{any} subsequent work evaluated the model on the scenario (\eg \citet{tay2022unifying} in the paper on \ultwo expanded the evaluation of \tfive to include \hellaswag and several other datasets) under \textbf{any} conditions (\eg fine-tuning, 0-shot prompting, 5-shot prompting).
After our evaluation (\textit{bottom}), models are now evaluated under the same conditions on many scenarios.
}

\paragraph{Standardization.}
To build a shared understanding of existing language models, consistent with our third element of holistic evaluation, we benchmark \nummodels prominent language models on \benchmarkname.
These models come from \nummodelcreators organizations: AI21 Labs (\eg \jurassicjumbo), Anthropic (\anthropic), BigScience (\eg \bloom), Cohere (\eg \coherexl), EleutherAI (\eg \gptneox), Google (\eg \ultwo), Meta (\eg \optonesevenfive), Microsoft/NVIDIA (\eg \mtnlgfivethreezero), OpenAI (\eg \gptdavinci), Tsinghua University (\glm), and Yandex (\yalm). 
Benchmarking these models is challenging given they vary in  accessibility \citep[see][]{liang2022community-norms}: some are open (\eg \gptneox), some are limited-access (\eg \gptdavinci), and some are closed (\eg \anthropic).
In some cases, very little is known about how these models were built (\eg the training data and its size are often not known), such as \instructdavinci. 
What we do know is that several of these models are deployed, either in external-facing commercial APIs (\eg the OpenAI playground) or products (\eg GitHub Copilot). 
That is, several of these models are having direct social impact at present. 
The absence of an evaluation standard compromises the community's ability to clearly and rigorously understand the overall landscape of language models.
To demonstrate how uneven language model evaluation has been, we annotated the datasets used to evaluate more than 40 language models (\ie all models evaluated in this work along with others like PaLM and Gopher) in \refapp{benchmark-comparison}.
We found major models such as \tfive and \anthropic were not evaluated on a single dataset in common in their original works \citep{raffel2019exploring, askell2021general}.
In fact, several models (\eg \jurassicgrande, \coherexl, \yalm) do not report any public results prior to our effort (to our knowledge).
And even for datasets that are frequently evaluated for across all \numlmdatasets datasets evaluated in major language modeling works (\eg \hellaswag; many of the datasets within GLUE and SuperGLUE), we find the evaluation conditions vary greatly.
On \hellaswag, some prior work reports fine-tuned accuracies (\eg \tfive), whereas others report prompting accuracies (\eg \gptdavinci).\footnote{
We emphasize our objective in raising this point is \textit{not} to suggest that any individual work has evaluated improperly. 
In fact, for this case, few-shot prompting was not even popularized at the time of T5's writing. 
But, nonetheless, these models are evaluated under very different conditions (\eg number of examples used in adaptation, ability to have white-box model access to use gradients to update the model), even if they are nominally evaluated on the same scenario.}
Even when works report results through few-shot prompting, the exact details can vary, which in \refsec{prompting-analysis} we show leads to wild swings in accuracies (\eg 30\% to 80\% for the same (model, scenario); see \citet{zhao2021calibrate}). 

In \reffig{evaluation-status-quo}, we make explicit how our evaluation changes the status quo.
Previously, on average models were evaluated on 17.9\% of our \core scenarios, even after compiling evaluations dispersed across different prior works.
We improve this to 96.0\%.\footnote{The remaining 4.0\% is due to technical issues with specific models that we document in \refsec{models}.}
By both evaluating these models on the same scenarios \emph{and} by conducting the evaluation under standardized conditions (\eg using the same few-shot prompting for all models), we facilitate direct head-to-head comparisons.

% Adaptation
\paragraph{The importance of adaptation.}
To benchmark these models, we must specify an \textit{adaptation} procedure that uses the general-purpose language model to tackle a given scenario \citep[see][\S4.3]{bommasani2021opportunities}.
In this work, we adapt all language models through few-shot prompting, as pioneered by GPT-3 \citep{brown2020gpt3}.
Furthermore, we opted to choose relatively simple, generic prompts in order to orient the development of language models towards generic language \emph{interfaces}
that respond robustly to direct natural language,
rather than requiring model-specific incantations.
Certainly stronger results could be obtained from
more sophisticated prompting \citep[\eg chain-of-thoughts;][]{wei2022chain},
prompt decomposition \citep{wu2022aichains, press2022selfask, arora2022ama}, 
and prompt-tuning \citep{lester2021power,li2021prefix},
potentially leading to qualitatively different findings \citep{suzgun2022challenging}.
The exploration of adaptation strategies is another dimension of benchmarking which we leave to future work.

\paragraph{Caveats and considerations.}
Before presenting our empirical findings, we highlight three key considerations.
First, while we standardize model evaluation, in particular by evaluating all models for the same scenarios, same metrics, and with the same prompts for \numincontext-shot prompting, models themselves may be more suitable for particular scenarios, particular metrics, and particular prompts/adaptation methods. 
To be explicit, while some models may perform poorly under our evaluation, they may perform well in other contexts.
Second, while the evaluation itself may be standardized, the computational resources required to train these models may be very different (\eg resource-intensive models generally fare better in our evaluation), which is partially captured by our measurements of efficiency.
Finally, models may also differ significantly in their exposure to the particular data distribution or evaluation instances we use, with the potential for \textit{train-test contamination}.
We emphasize that we have a limited understanding on how contaminated models are, and to what extent this compromises the validity and legitimacy of our evaluation, though we do provide all evidence we are aware of in \refapp{contamination}. 

\hypertarget{findings}{\subsection{Empirical findings}}
\label{sec:findings}
To give a sense of the magnitude of our evaluation, we ran a total of \numrunresults runs (\ie evaluating a specific model on a specific scenario).
% , which are all available at \benchmarkUrlRuns.
This amounts to a total cost of \totaltokens tokens and 
\totalqueries queries across all models, 
\$\totaldollars for the commercial APIs, and about \totalgpuhours \gpuhours worth of compute for the open models. \\

\noindent Here is a summary of the high-level findings:
\begin{enumerate}
    \item \textbf{The benefits of instruction-tuning.}
    Across the \core scenarios, we find that \instructdavinci performs best on our accuracy, robustness, and fairness metrics, with \anthropic being in the top 3 for all 3 metrics (despite being more than 10$\times$ smaller in model scale compared to \mtnlgfivethreezero, which is the second most accurate and fair) as shown in \reffig{head-to-head}. 
    Given the very strong performance of both models, and that they are the only instruction-tuned models we evaluate (beyond the much smaller OpenAI model variants), this suggests instruction-tuning provides a broad set of advantages.

    \item \textbf{Relating model accuracy with model access.}
    In light of the high accuracies of \anthropic (closed), \mtnlgfivethreezero (closed), and \instructdavinci (limited-access), we observe a consistent gap on all \core scenarios (\reffig{accuracy-vs-access}) between the current open models and non-open models.
    We emphasize that this gap reflects the current snapshot of models we evaluate (\reftab{models}), and that the gap could grow or shrink over time as new models are released.
    On one hand, we see the recent release of open models (\optonesevenfive, \bloom, \glm) as greatly reducing the gap over the past year, but we also have not evaluated some non-open models (\eg PaLM, Gopher) that we expect to be quite accurate.
    In either case, monitoring this gap over time is crucial for 
    tracking the accessibility (or lack thereof) and ultimately the power dynamics associated with language models.
    
    \item \textbf{Calibration.}
    We observe that the relationship between accuracy and calibration (\refsec{metrics-calibration}) depends on the scenario and adaptation procedure (\reffig{acc-X}, \reffig{inter-metric-correlations}).
    As an example, for \hellaswag,\footnote{See \benchmarkUrlPrefix{hellaswag}.} improving accuracy worsens calibration,
    whereas for \openbookqa,\footnote{See \benchmarkUrlPrefix{openbookqa}.} improving accuracy improves calibration.
    
    \item \textbf{Robustness and fairness perturbations.}
    Across all scenarios, we observe strong correlations between accuracy, robustness, and fairness, where robustness and fairness metrics consider worst-case accuracy over a set of perturbations (\eg typos for robustness, dialect alteration for fairness)---see \refsecs{metrics-robustness}{metrics-fairness} for more details.
     While there is a strong correlation between accuracy and fairness (\reffig{acc-X}, \reffig{inter-metric-correlations}), we do observe trade-offs where the most accurate model is not the most robust or most fair.
    We also see serious drops in some cases: for example, on \narrativeqa, \mtnlgfivethreezero precipitously drops from 72.6\% standard accuracy (\ie the third-most accurate model) to 38.9\% accuracy in the presence of robustness perturbations.\footnote{See \benchmarkUrlPrefix{narrative_qa}.}
    
    \item \textbf{Performance disparities.}
    When we have access to demographic metadata, we generally see consistent performance disparities for all models. 
    As an example of racialized dialect disparities, \optonesevenfive is the most accurate model on \twitteraae but its accuracy degrades from 1.506 bits per byte for White English to 2.114 bits per byte for African American English (lower is better).\footnote{See \benchmarkUrlPrefix{twitter_aae}.}
    
    \item \textbf{Generative harms.}
    We find that the biases and toxicity in model generations are largely constant across models and low overall on average for the \core scenarios (\reffig{acc-X}).
    However, note that even low levels of bias or toxicity could cause non-trivial social harm, and targeted evaluations are needed to obtain a more detailed characterization (\refsecs{targeted-bias}{targeted-toxicity}).
    
    \item \textbf{Accuracy vs. efficiency.}
    We do not see a strong trade-off between accuracy and efficiency (which depends on both the model architecture and the hardware, see \refsec{metrics-efficiency}) across all \nummodels models (\reffig{acc-X}).
    For each family of models (\eg different size variants of GPT-3), we find that as models become larger, accuracy consistently improves but with higher training and inference cost.\footnote{See \benchmarkUrlPrefix{core_scenarios\#Efficiency}.}
    Overall, we observe that only a subset of all models (across model families) are on the accuracy-efficiency Pareto frontier for each scenario.

    \item \textbf{Question answering.}
    Across the 9 \core question answering scenarios (\refsec{questionAnswering}), we observe significant heterogeneity in results, though \instructdavinci is the most accurate model for all 9 scenarios.\footnote{See \benchmarkUrlPrefix{question_answering}.}
    In fact, for 6 of the 9 scenarios, there is no open model among the three most accurate models, as generally they are \instructdavinci, \anthropic, and \mtnlgfivethreezero in descending order of accuracy.
    
    \item \textbf{Information retrieval.}
    We consider the classic task of ranking candidate passages given a query (\refsec{informationRetrieval}).
    The best-performing models we evaluate outperform classical retrieval methods and under some settings perform comparably to various fine-tuned neural retrievers, while nonetheless trailing the state of the art.\footnote{See \benchmarkUrlPrefix{information_retrieval}.}
    Because the number of candidates could be large,
    we create a LM request per passage, which requires
    the model to produce calibrated probabilities.
    Our use of LMs for passage ranking is unorthodox,
    and computationally intensive in its naive implementation,
    but we include it as a proof of concept.
    %These promising results highlight some key challenges in ranking tasks. Information retrieval (IR) requires sorting documents by relevance to a query, demanding calibrated assessments of relevance \textit{across} documents.
    %Naively, IR also requires prompting language models numerous times per query (\ie for each document), which---at current latencies---would be sufficiently fast for typical applications only whenever parallelization across documents is cost-effective.
    
    \item \textbf{Summarization.}
    \cnndm and \xsum have been standard benchmarks for summarization for many years, but we find that automated evaluations on these datasets largely fail to discriminate differences we observed in model quality.
    % Removed for TMLR
    % the official reference summaries in these datasets are outperformed by generated model summaries in human evaluations, especially for faithfulness (\reftab{summarization-human-evaluation}).
    % Overall, we believe summarization benchmarks (along with metrics) must be improved by incorporating high-quality, human-written summaries (\refsec{missing-scenarios}), so that we can draw meaningful conclusions on the effect of in-context learning, instruction tuning, and fine-tuning (see \refsec{human-evaluation-summarization}).
    
    \item \textbf{Sentiment analysis.}
    For sentiment analysis on \imdb, many models are quite accurate and well-calibrated with marginal drops on robustness and fairness perturbations, but the contrast sets of \citet{gardner2020contrast} highlight clear limitations in model robustness (\eg one of the most accurate models in \glm drops by more than 8\%).\footnote{See \benchmarkUrlPrefix{sentiment_analysis} and \benchmarkUrlPrefix{robustness_contrast_sets}.} 
    
    \item \textbf{Toxicity detection.}
    For toxicity detection on \civilcomments, we find that most models are not particularly accurate: \optonesevenfive is one of the most accurate models across all scenarios (\reffig{head-to-head}), but achieves essentially chance accuracy at 50.1\%.\footnote{See \benchmarkUrlPrefix{civil_comments}.}
    Critically, given the importance of fairness in toxicity detection due the disparate impacts of content moderation, we find that most models are similarly accurate for detecting toxicity in comments mentioning Black and White individuals.
    However, models vary greatly in their robustness: \optonesevenfive drops from 51.3\% standard accuracy to 8.8\% robust accuracy on the Black split, whereas the drop is less precipitous on the White split (50.8\% to 24.3\%). 
    
    \item \textbf{Miscellaneous text classification.}
    For text classification on \raft, we see significant heterogeneity in which models do well on which subsets/tasks.\footnote{See \benchmarkUrlPrefix{raft}.}
    \instructdavinci is consistently accurate across splits when compared with other models, but performs very poorly on the Systematic Review Inclusion split with an accuracy of 40.8\% compared to 97.5\% from several models (\eg \glm). 
    
    \item \textbf{Linguistic understanding.}
    The trends in accuracy for language modeling\footnote{See \benchmarkUrlPrefix{language\#Accuracy}.} are quite different from the trends for the \core scenarios (\reffig{head-to-head}) .
    In particular, \gptneox, \optonesevenfive, \bloom, \gptj, and \optsixsix consistently have the lowest bits-per-byte (lower is better) on \pile, \twitteraae, and \ice. 
    In terms of linguistic phenomena, all models perform fairly similarly on \blimp overall, and further perform very similarly even on each of the specific subsets for morphology, syntax, semantics, and syntax-semantics.
    We see the widest spread on irregular forms (morphology), where surprisingly the models that tend to be the most accurate for core scenarios (\ie \instructdavinci, \mtnlgfivethreezero) are some of the least accurate for irregular forms, perhaps suggesting they have overgeneralized particular linguistic rules.\footnote{See \benchmarkUrlPrefix{blimp\#phenomenon:\%20irregular_forms}.}
    
    \item \textbf{Knowledge.}
    \instructdavinci demonstrates superior performance for all knowledge-intensive evaluations,\footnote{See \benchmarkUrlPrefix{knowledge\#Accuracy}.} with a very sizable gap for accuracy on \truthfulqa of 62.0\% compared to second place of 36.2\% from \anthropic.\footnote{See \benchmarkUrlPrefix{truthful_qa}.    
    We note this is especially interesting given the projections of model accuracy by \citet{evans2022truthfulqa}, though we note our results are for \numincontext-shot learning whereas their projections are for 0-shot learning.}
    Further, \mtnlgfivethreezero shows strong performance on the highly knowledge-intensive \naturalquestions (closed-book) and \wikifact scenarios, which generally concurs with the hypothesis that model scale especially contributes to improvements in acquisition of factual knowledge.
    For example, \anthropic and \mtnlgfivethreezero tend to get very similar accuracies for most scenarios (as suggested by \reffig{head-to-head}), but \mtnlgfivethreezero demonstrates a wide margin for these two scenarios (38.5\% vs. 28.7\% for \naturalquestions (closed-book), 34.3\% vs. 22.3\% for \wikifact). 
    
    \item \textbf{Reasoning.}
    For reasoning-intensive scenarios, we find that the code models, especially \davincicodex, consistently outperform the text models, even on synthetic reasoning scenarios posed in natural language.\footnote{See \benchmarkUrlPrefix{reasoning\#Accuracy}.}
    This gap is made clear in mathematical reasoning: for \gsm, \davincicodex achieves an accuracy of 52.1\%, where the next best model is \instructdavinci at 35.0\% and no other model surpasses 16\%.\footnote{See \benchmarkUrlPrefix{gsm}.} 
    Further, in addition to \davincicodex, \instructdavinci is much more accurate than other text models (\eg 65.1\% accuracy on synthetic reasoning in natural language, whereas the next most accurate text model is \optonesevenfive at 29.4\% accuracy, and \davincicodex has an accuracy of 72.7\%). 
    
    \item \textbf{Memorization of copyrighted/licensed material.}
    We find that the likelihood of direct regurgitation of long copyrighted sequences is somewhat uncommon, but it does become noticeable when looking at popular books.\footnote{See \benchmarkUrlPrefix{copyright_text}.}
    However, we do find the regurgitation risk clearly correlates with model accuracy: \instructdavinci, \gptdavinci, and \anthropic demonstrate the highest amount of verbatim regurgitation in line with their high accuracies. 
    
    \item \textbf{Disinformation.}
    We find that the largest models (particularly \instructdavinci and \anthropic) are effective at generating realistic headlines that support a given thesis,\footnote{See \benchmarkUrlPrefix{disinformation_reiteration}.} but results are more mixed when prompting models to generate text encouraging people to perform certain actions (\reftab{disinformation-human-evaluation}).\footnote{See \benchmarkUrlPrefix{disinformation_wedging}.}
    
    \item \textbf{Targeted biases.}
    For \bbq, \instructdavinci is the most accurate model by a very wide margin (89.5\% accuracy), with the next most accurate models (\tzero, 48.4\%; \mtnlgfivethreezero, 44.9\%) being the only other models with accuracies above 40\%. 
    We highlight this because we see a very striking relationship on \bbq between model accuracy and model bias for ambiguous contexts.
    These three models, which are the three most accurate, are the only three models with biases in ambiguous contexts that align with broader social biases/discrimination, whereas all other models show biases in the other direction (\reffig{component-bias}).
    In other words, we find that for \bbq the most accurate models are precisely those that are most concerning for social biases in ambiguous contexts, though the trends in disambiguated contexts are less clear.
    
    \item \textbf{Targeted toxicity generation.}
    For the \core scenarios, we observed the rate of toxicity generation was quite low.
    Honing in on toxicity generation, all models show much stronger tendencies for toxic generations for toxic prompts in \realtoxicityprompts, as compared to relatively non-toxic prompts in both \realtoxicityprompts and \bolddataset.\footnote{See \benchmarkUrlPrefix{harms\#Toxicity}.} 
    Understanding how these trends change based on the automated toxicity detection model used (currently PerspectiveAPI), as well as when human judgments from diverse stakeholders are used, is a key area for future work.
    
    \item \textbf{Comprehensiveness.}
    By evaluating under unified conditions extensively, we expose findings lying in plain sight. 
    In other words, while in many cases we are evaluating models that are available publicly on datasets that are available publicly, we nonetheless surface new findings. 
    As an example, we find \instructdavinci achieves an accuracy of 74.4\% \texttt{ROUGE-L} on \narrativeqa, which sets a new state-of-the-art across all methods to our knowledge, in this case over the strong QA-specialized  \textsc{UnifiedQA}-v2 model \citep[67.4\% \texttt{ROUGE-L};][]{khashabi2022unified}. 
    
    \item \textbf{Prompting.}
    All models show significant sensitivity to the formatting of prompt, the particular choice of in-context examples, and the number of in-context examples across all scenarios and for all metrics (see \refsec{prompting-analysis}).
    In this effort, we consistently work towards standardizing these dimensions (\eg to ensure models are interoperable/performant using the same prompting practices), but current models differ in what prompting decisions would maximize accuracy.\footnote{See \benchmarkUrlPrefix{ablation\_prompts}.}

    \item \textbf{Multiple choice adaptation method.}
    We find that model performance is extremely sensitive to how multiple choice scenarios are adapted into prompts: for example, accuracy for \optonesevenfive on \hellaswag is 79.1\% when each answer choice is presented in a separate 0-shot prompt (\ie one of the most accurate models), but drops precipitously to 30.2\% (almost random accuracy) when the answer choices are presented jointly in a single \numincontext-shot prompt (\ie in the format of a multiple-choice exam).\footnote{See \benchmarkUrlPrefix{ablation_multiple_choice}.} 
    Further, even for the same scenario, the adaptation method that maximizes accuracy can differ (and produce qualitatively different results) across models (\reffig{mc-ablations}) .
    This poses a fundamental challenge for what it means to standardize language model evaluation in a fair way across models.
    
    \item \textbf{Upstream perplexity and downstream accuracy.}
    Given the myriad scenarios where LMs could provide value, it would be appealing for many reasons if upstream perplexity on language modeling objectives reliably predicted downstream accuracy.
    Unfortunately, when making these comparisons across model families, even when using bits-per-byte (BPB; which is more comparable than perplexity), we find this type of prediction does not work well: BPB on \pile is a poor predictor of downstream accuracy (\reffig{perplexity-accuracy}) though we note some models are trained on \pile whereas others are not (\reftab{contamination-evidence}).
    More broadly, given the many downstream results, we encourage future work to explore new intrinsic/upstream surrogate measures of performance that can be shown to reliably predict downstream results (including for desiderata beyond accuracy) as discussed in \citet[][\S4.4.2]{bommasani2021opportunities} 
    
    \item \textbf{Trends for model scale.}
    We find that model scale, within a model family, reliably predicts model accuracy, but for no scenario is a good predictor of downstream accuracy across all models (\reffig{scale-accuracy}).
    However, we see a very clear thresholding effect: all models that win head-to-head model comparisons for accuracy at a rate well above chance (\ie $>55\%$) are at least 50B parameters (\reffig{head-to-head}).
    Of these models, which are the 10 most accurate models, some of the most accurate (\ie in the top 5) are the smallest (\anthropic, \coherexl).
    Overall, scale seems to be a key determinant of accuracy, and scaling within a model family reliable improves accuracy, but it might be inefficient compared to other means (\eg training with human feedback; compare \mtnlgfivethreezero and \anthropic).
\end{enumerate}
\hypertarget{contributions}{\subsection{Contributions}}
\label{sec:contributions}
To summarize, our contributions are:
\begin{enumerate}
  \item \textbf{Taxonomy.}
  We taxonomize the vast design space of language model evaluation into scenarios and metrics.
  By stating this taxonomy, we can select systematically from this space, which makes explicit both \textit{our priorities} in benchmark design and \textit{the limitations} in the benchmark at present (see \refsec{missing}).
  \item 
  \textbf{Broad coverage.}
  Given our taxonomy, we select and implement \numcorescenarios \core scenarios, for which we comprehensively measure \nummetriccategories metrics (accuracy, calibration, robustness, fairness, bias, toxicity, efficiency).
  We also include \numcomponents targeted evaluations of skills and risks (\eg knowledge, reasoning, disinformation, copyright), introducing \numnewscenarios new scenarios that have not been previously used in mainstream language model evaluation.
  \item \textbf{Evaluation of existing models.}
  We evaluate \nummodels language models under the standardized conditions of our benchmark, ensuring models can now be directly compared across many scenarios and metrics.
  These models vary in terms of their public accessibility: \numpublicmodels are open, \numlimitedmodels are limited-access, and \numclosedmodels are closed.
 \item \textbf{Empirical findings.}
 Our extensive evaluation yields a host of findings (\refsec{experiments}), which in some cases reinforce findings in the literature and in others produce new knowledge about today's language models.
 These results offer guidance for future language model development and ample opportunities for further analysis. 
  \item \textbf{Interactive results and codebase.}
  We provide a public website with all results, underlying model predictions and adaptation details, along an extensible codebase to support the community in taking \benchmarkname further.\footnote{\benchmarkUrl{}}
\end{enumerate}

\paragraph{Acknowledging the prior work this effort builds on.}
To build our holistic evaluation of language models, we directly build on top of many prior works.
While we advocate for evaluating language models in their totality, \ie centralizing many disparate evaluations, we want to be explicit that the underlying works across the AI community \textbf{should be recognized and cited}, as \benchmarkname would not exist in its current form without them.
In particular, if the results of \benchmarkname are used by future work or new models are evaluated on \benchmarkname, they should cite the works that created the many datasets/evaluations that constitute \benchmarkname.\footnote{We take direct inspiration from and follow the precedent set by GLUE and SuperGLUE \citep{wang2019superglue, wang2019glue}; see \url{https://super.gluebenchmark.com/faq}.} 
For this reason, we provide the BibTeX entries for all of these works in the codebase\footnote{\githubUrl{}} and explicitly acknowledge the associated work for every evaluation on the website.\footnote{\benchmarkUrl{}}